\paragraph{Baseline (base paper, English)}
\begin{small}\begin{verbatim}
Please analyze the image carefully, reason step-by-step and provide
your answer at the end.

For this question:
- Provide a numerical value
- Round to appropriate decimal places if needed
- Format your answer in \boxed{} (e.g., \boxed{2.5} or \boxed{42})
\end{verbatim}\end{small}

\paragraph{PoT solver instructions}
\begin{small}\begin{verbatim}
How to solve:
- Write the solution as numbered steps, each starting on a new line
with "Step 1:", "Step 2:", and so on.
- Explain the physics/chemistry/mathematics of each step in words.
- Do NOT do calculations in your head. Whenever a step needs a
calculation (arithmetic, algebra, solving equations, derivatives,
integrals, limits, unit conversions, or checking each option), write a
short Python program in a ```python code block that prints the result.
You may use math, fractions, sympy and numpy.
- End the code block with ```. The code will be run and its printed
output will be given back to you in an ```output block. Then continue
the solution using that output. Never invent the output yourself.
- Use at most 3 code blocks.
- Be concise: do not re-check the same step again and again.
- Finish with exactly one final answer in \boxed{}. Do not use \boxed
anywhere else. For a numerical answer write a plain decimal number
inside the box (no units, no fractions, no LaTeX), rounded to two
decimal places if it is not an integer.

Required format (example):
Step 1: <which law / concept applies and the equation to use>
```python
import sympy as sp
x = sp.symbols('x')
print(sp.solve(2*x - 6, x))
```
```output
[3]
```
Step 2: <interpret the printed result, check the options>
\boxed{<answer>}
\end{verbatim}\end{small}

\paragraph{Critic}
\begin{small}\begin{verbatim}
You are a strict examiner checking a student's solution to a JEE
Advanced question. The question is in the attached image{lang_note}.

Instructions the student was given:
{type_instruction}

Student's solution (steps are numbered):
<solution>
{solution}
</solution>

Check the steps in order: reading of the question and figure, chosen
concept and equations, algebra and arithmetic (the ```output blocks
are real program outputs), and the final option or number. Find the
FIRST step that contains an error. Do not solve the whole problem
again. Keep your check short (at most about 250 words), then end your
reply with exactly these three lines:
VERDICT: CORRECT or ERROR
FIRST_ERROR_STEP: <step number, or NONE>
EXPLANATION: <one or two sentences: what is wrong and how to fix it>
\end{verbatim}\end{small}

\paragraph{Corrector feedback}
\begin{small}\begin{verbatim}
A reviewer checked an earlier attempt at this question. The steps
before Step {k} were accepted and are already written at the start of
your answer. The reviewer found an error in Step {k}:
"{explanation}"
Continue the solution from Step {k}: re-derive that step correctly (do
not repeat the error) and complete the solution, following the same
rules (Python code for every calculation, one final \boxed{{}}
answer).
\end{verbatim}\end{small}

\paragraph{Concept tagger}
\begin{small}\begin{verbatim}
Look at this JEE Advanced exam question (it may be written in English
or Hindi). Do not solve it. In English, describe what is needed to
solve it, using exactly this format on one line:
SUBJECT: <Physics/Chemistry/Mathematics> | TOPIC: <specific topic, 2-6
words> | METHOD: <key concept or solution technique, at most 15 words>
\end{verbatim}\end{small}

\paragraph{Transcriber}
\begin{small}\begin{verbatim}
Transcribe this exam question exactly as plain text in its original
language. Write mathematical expressions in LaTeX. Include all answer
options. If there is a figure, add one line "[Figure: <short
description>]". Output only the transcription.
\end{verbatim}\end{small}

\paragraph{Exemplar header}
\begin{small}\begin{verbatim}
Below are worked solutions of earlier JEE Advanced problems that use a
similar method. Use them only as guidance on the method. They are
DIFFERENT problems: do not copy their numbers or answers.
\end{verbatim}\end{small}
